\chapter{Implementation}
\label{ch:implementation}

\section{Implementation Goals}
The delivered artefact is a dependency-free browser application designed for interactive exploration of distributed mutual exclusion. The implementation prioritises correctness (especially the safety invariant), pedagogical clarity, and reproducibility. Two representative approaches are implemented:
\begin{itemize}
  \item a token-based \textbf{Token Ring} model;
  \item a message-based \textbf{Ricart--Agrawala (RA)} teaching-oriented prototype \cite{ricart1981optimal,lamport1978time}.
\end{itemize}

\section{Repository Structure and Responsibilities}
The codebase is organised around a small number of responsibilities:
\begin{itemize}
  \item \texttt{index.html}: page layout and controls.
  \item \texttt{styles.css}: minimal CSS for layout and readability.
  \item \texttt{mutex\_core.js}: algorithm models, state transitions, safety checking, script replay, import/export.
  \item \texttt{mutex\_main.js}: DOM wiring, rendering, preview drawing, evidence export, and interaction logic.
  \item \texttt{examples/}: scripted JSON scenarios.
  \item \texttt{tools/smoke\_tests.mjs}: dependency-free automated smoke tests.
\end{itemize}

This separation keeps the simulation engine testable without a browser and prevents UI changes from implicitly changing algorithm behaviour \cite{gamma1994design}.

\section{Core Data Structures}
A single \texttt{model} object captures the current simulator state. It stores:
\begin{itemize}
  \item the selected algorithm and current mode,
  \item the list of processes,
  \item algorithm-specific state (token information or message queue),
  \item a step-numbered trace,
  \item and lightweight metrics.
\end{itemize}

\subsection{Process State}
Each process includes:
\begin{itemize}
  \item \texttt{id}
  \item \texttt{requesting}
  \item \texttt{inCS}
  \item \texttt{crashed}
\end{itemize}

For RA, additional fields include:
\begin{itemize}
  \item \texttt{clock}
  \item \texttt{reqTs}
  \item \texttt{awaiting}
  \item \texttt{deferred}
\end{itemize}

\subsection{Token Ring State}
Token Ring uses:
\begin{itemize}
  \item \texttt{ring}
  \item \texttt{token.holder}
  \item \texttt{token.lost}
\end{itemize}

\subsection{RA Network State}
The RA prototype uses:
\begin{itemize}
  \item \texttt{network.queue}
  \item \texttt{network.nextMsgId}
  \item \texttt{network.dropNextSend}
\end{itemize}

\subsection{Trace and Metrics}
A step-numbered trace records user actions, algorithm events, warnings, and stall explanations. Metrics record token passes, CS entries, releases, and for RA the number of messages sent, delivered, and dropped.

\section{Safety Checking}
The safety property is encoded directly in the core module through a dedicated checker. The invariant is:
\begin{quote}
At most one process may be in the critical section at any time.
\end{quote}
The UI exposes this as a Safety indicator, making the invariant visible during every scenario.

\section{Step Engine}
A core design feature is the \textbf{step engine}. Each press of \texttt{Step} performs one minimal unit of progress.

\subsection{Token Ring Step Semantics}
A Token Ring step follows this decision order:
\begin{enumerate}
  \item If the token is lost, report that progress cannot continue.
  \item If a process is already in the critical section, report that release is required.
  \item If the token holder is requesting, allow entry to the critical section.
  \item Otherwise, pass the token to the next process in ring order.
\end{enumerate}

\subsection{RA Step Semantics}
An RA step follows this decision order:
\begin{enumerate}
  \item If the network queue is non-empty, deliver exactly one queued message.
  \item If the queue is empty but a process is still awaiting replies, emit a stall explanation.
  \item Otherwise, report that no messages are in flight.
\end{enumerate}

\section{Token Ring Implementation}
Token Ring provides mutual exclusion through token ownership. A process requests the critical section by setting \texttt{requesting = true}. If the token holder is requesting and no process is already in the critical section, it enters.

\subsection{Token Faults and Recovery}
Token Ring supports:
\begin{itemize}
  \item \textbf{Drop token}: sets \texttt{token.lost = true}, halting progress.
  \item \textbf{Regenerate token}: restores the token.
  \item \textbf{Crash / recover}: makes a process unavailable and later returns it to participation.
\end{itemize}

\section{Ricart--Agrawala Implementation}
In the RA prototype, a requesting process:
\begin{enumerate}
  \item increments its logical clock,
  \item stores the request timestamp,
  \item initialises its outstanding reply set,
  \item and enqueues REQUEST messages to all peers.
\end{enumerate}

When a REQUEST is delivered, the receiver either:
\begin{itemize}
  \item replies immediately, or
  \item defers reply if it is already in the critical section or has priority.
\end{itemize}

Entry occurs only after all required REPLY messages are received.

\subsection{Tie-break Logic}
When requests share the same timestamp, the implementation resolves the tie using numeric process IDs, ensuring deterministic and reproducible behaviour.

\section{Fault Injection Points}
The fault controls are implemented in the simulation engine.

\subsection{Token Ring}
\begin{itemize}
  \item \texttt{dropToken(model)}
  \item \texttt{regenerateToken(model)}
  \item \texttt{crashProcess(model, pid)}
  \item \texttt{recoverProcess(model, pid)}
\end{itemize}

\subsection{Ricart--Agrawala}
\begin{itemize}
  \item \texttt{dropNextMessage(model)}
  \item \texttt{toggleDropNextSend(model)}
  \item crash/recover controls
\end{itemize}

\section{Stall Explanation Logic}
A key clarity feature is explicit stalled-state explanation. In RA, if the queue is empty but a process still awaits replies, the simulator emits a trace message indicating which REPLY is missing. This prevents learners from misinterpreting ``queue empty'' as successful completion.

\section{Scripted Scenarios}
Scripted scenarios are stored as JSON files and replayed via the same step engine used in interactive mode. This supports repeatable demonstrations and reproducible evidence capture.

\section{Export Pipeline}
The tool exports:
\begin{itemize}
  \item state JSON,
  \item trace TXT,
  \item preview PNG.
\end{itemize}

The export naming strategy combines a timestamp-based session identifier, algorithm/mode metadata, and a sequence number, producing stable and traceable evidence artefacts.

\section{Implementation Limitations}
The implementation intentionally simplifies some aspects:
\begin{itemize}
  \item the RA network is modelled as a FIFO queue,
  \item no retransmission or failure detector is implemented,
  \item the visual design prioritises clarity over animation.
\end{itemize}

\IfFileExists{figures/fig_ui_overview.png}{
\begin{figure}[htbp]
\centering
\includegraphics[width=0.95\linewidth]{fig_ui_overview.png}
\caption{Main user interface of the Distributed Mutual Exclusion Explorer.}
\label{fig:ui-overview}
\end{figure}
}{}